\section{引言：什么是 LLM Agent？}

Shunyu Yao（姚顺雨）是 ReAct 框架的提出者，现就职于 OpenAI。本次讲座系统梳理了 LLM Agent 的定义、历史发展脉络和未来方向。

\subsection{Agent 的定义}

\begin{importantbox}{Agent 的三层定义}
\begin{enumerate}
    \item \textbf{Text Agent}：与环境交互的智能系统，其动作（action）和观测（observation）均以语言（text）形式表达。Text Agent 不一定使用 LLM，早在 AI 诞生之初就有基于规则的 Text Agent。
    \item \textbf{LM Agent}：使用语言模型来生成动作的 Text Agent。
    \item \textbf{Reasoning Agent}：使用语言模型\textbf{推理后再行动}的 Agent——这是一个质的飞跃。
\end{enumerate}
\end{importantbox}

定义 Agent 需要回答两个子问题：什么是``智能''（intelligence）？什么是``环境''（environment）？

\begin{itemize}
    \item 环境可以是物理的（机器人、自动驾驶）或数字的（视频游戏、网页、代码编辑器）。
    \item ``智能''的定义随时代演变——60 年前简单聊天机器人就算智能，如今 ChatGPT 都不再令人惊讶。
\end{itemize}

\subsection{本章小结}

LLM Agent 的核心创新不是让语言模型执行动作，而是让它\textbf{先推理再行动}。Reasoning Agent 区别于此前所有 Agent 范式的关键在于：推理（thinking）是一种内部动作，具有无限的语言空间。

